% ============================================================================
%  Compiles with pdfLaTeX. Diagrams are ASCII inside fancyvrb.
% ============================================================================
\documentclass[11pt,a4paper]{article}

\usepackage[T1]{fontenc}
\usepackage[utf8]{inputenc}
\usepackage[margin=2.2cm]{geometry}
\usepackage{amsmath,amssymb}
\usepackage{booktabs}
\usepackage{tabularx}
\usepackage{array}
\usepackage{fancyvrb}
\usepackage{enumitem}
\usepackage[hidelinks]{hyperref}

\newcolumntype{L}{>{\raggedright\arraybackslash}X}
\setlength{\tabcolsep}{4pt}
\renewcommand{\arraystretch}{1.15}

\DefineVerbatimEnvironment{code}{Verbatim}{fontsize=\small,xleftmargin=1em}
\DefineVerbatimEnvironment{codetiny}{Verbatim}{fontsize=\footnotesize,xleftmargin=1em}

\newcommand{\yhat}{\hat{y}}
\newcommand{\ind}{\mathbf{1}}

\title{Three-dataset probabilistic baselines: data and models}
\author{}
\date{}

\begin{document}
\maketitle
\vspace{-3em}

\noindent
Reference card for \texttt{elia} / \texttt{ercot} / \texttt{spain}. Two models are built and
run: \textbf{A} density combination (Li, Wang \& Zhang 2020) and \textbf{B} iQRNN (Zhang, Quan
\& Srinivasan 2019). QRA (Liu et al.\ 2017) was dropped --- its sister family is driven by a
temperature recency effect and none of the three datasets has temperature.

Code: \texttt{src/bench/}, entries \texttt{scripts/train\_bench\_\{density,iqrnn\}.py},
results \texttt{results/bench/}.

% ============================================================================
\section{Datasets}

\begin{center}
\small
\begin{tabular}{@{}l r r r@{}}
\toprule
 & \texttt{elia} & \texttt{ercot} & \texttt{spain} \\
\midrule
granularity                & 15 min & 60 min & 60 min \\
steps per day $H$          & 96 & 24 & 24 \\
lookback $L$ (steps)       & 672 & 168 & 168 \\
deadline                   & $(D{-}1)$ 23:00 & $(D{-}1)$ \textbf{09:00} & $(D{-}1)$ 23:00 \\
deadline gap $G$ (steps)   & 4 & 15 & 1 \\
\textbf{lag window}        & $[99,\,675]$ & $[38,\,182]$ & $[24,\,168]$ \\
days train/val/test        & 2239 / 748 / 746 & 416 / 142 / 140 & 869 / 292 / 291 \\
steps train                & 214{,}944 & 9{,}984 & 20{,}856 \\
target                     & \texttt{Total Load} & \texttt{ACTUAL\_NetLoad} & \texttt{total load actual} \\
\bottomrule
\end{tabular}
\end{center}

\noindent
A lag $k$ is usable for \emph{every} step of the target day iff
\begin{equation}
  G + H - 1 \;\leq\; k \;\leq\; L - 1 + G .
  \label{eq:window}
\end{equation}
Lower bound: the day's last step is $H{-}1$ past its first, itself $G$ past the deadline.
Upper bound: anything deeper reaches past the start of the lookback window. Every feature list
below is filtered through this.

\subsection{Model A features --- target lags only}

\texttt{10--11 target lags $+$ hour $+$ dayofweek $+$ month $+$ is\_weekend.} No exogenous
columns at all: Li's design is that the base learners are generic, so the combination step is
what is tested.

\medskip
\noindent
Published list (hours): $24,25,26,27,28,48,72,96,120,144,168$. Instantiated per dataset by
shifting \emph{whole days} until the smallest member clears Eq.~\eqref{eq:window}'s lower
bound, then dropping anything past the upper bound.

\begin{tabularx}{\textwidth}{@{}l c p{5.2cm} L@{}}
\toprule
 & $n$ & lag list (steps) & note \\
\midrule
\texttt{spain} & 11 & \texttt{24,25,26,27,28,48,72,\allowbreak 96,120,144,168}
  & the paper's list verbatim \\
\texttt{ercot} & 10 & \texttt{48,49,50,51,52,72,96,\allowbreak 120,144,168}
  & shifted $+24$ (deadline 9, so $D{-}1$ is not fully observed); 192 exceeds 182, dropped \\
\texttt{elia}  & 10 & \texttt{192,196,200,204,208,\allowbreak 288,384,480,576,672}
  & shifted $+96$; 768 exceeds 675, dropped \\
\bottomrule
\end{tabularx}

\medskip
\noindent
So \texttt{ercot} and \texttt{elia} carry 14 features, \texttt{spain} 15. Both losses are
forced: no shift satisfies both ends with eleven lags.

\subsection{Model B features --- pool $\rightarrow$ RFE $\rightarrow$ 103}

The paper's pool is calendar $+$ load lags $+$ \textbf{temperature}. No dataset has
temperature, so it is replaced by whatever exogenous columns each provides.

\begin{center}
\small
\begin{tabular}{@{}l l l l l r@{}}
\toprule
 & target lags & \texttt{past\_cols} & \texttt{known\_cols} & calendar & pool \\
\midrule
\texttt{elia}  & 3{,}262 (99\ldots3360, every step) & --- & $4 \times$ lag 0 & 11 & 3{,}277 \\
\texttt{ercot} & 803 (38\ldots840)  & $11 \times 8$ echoes $= 88$  & --- & 11 & 902 \\
\texttt{spain} & 817 (24\ldots840)  & $22 \times 7$ echoes $= 154$ & $1 \times$ lag 0 & 11 & 983 \\
\bottomrule
\end{tabular}
\end{center}

\begin{itemize}[leftmargin=1.4em,itemsep=1pt]
  \item \textbf{target}: dense grid, floor $\rightarrow$ 5 weeks, every step.
  \item \textbf{\texttt{past\_cols}}: daily-echo grid
        $\{\mathrm{floor}\} \cup \{\lceil \mathrm{floor}/H \rceil H + kH,\ k = 0\ldots6\}$.
        A dense grid on \texttt{spain}'s 22 series would be an 18{,}000-column pool.
  \item \textbf{\texttt{known\_cols}}: lag 0 only (day-ahead forecast, published pre-deadline).
  \item \textbf{calendar}: HoD / DoW / MoY raw, three sine--cosine pairs, \texttt{is\_weekend}.
\end{itemize}

\noindent
RFE (random forest, MAE criterion, \texttt{step=0.35}) on a 12{,}000-row training subsample
cuts the pool to $100 + 3$ manual add-backs $= 103$.

\medskip
\noindent
\textbf{Caveat.} The \texttt{known\_cols} are misaligned in the source files --- they correlate
best with the target at $+24$ steps, not lag 0 (\texttt{elia}'s two forecast columns at $+24$
and $+168$, which are the \emph{hourly} constants applied to a 15-minute series). Used
unshifted as configured. Degradation, not leakage; hits \texttt{ercot} hardest since its target
is net load and its wind/solar forecasts are the main signal. Model A is unaffected.

% ============================================================================
\section{Models}

\subsection{Model A --- density combination}

\begin{codetiny}
                    T1 = train        T2 = val 1st half
                    T3 = val 2nd half T4 = test

   X (n_steps, 14 or 15)
        |
        +--> GPR  x5   ARD Matern3/2 + RBF + White; 2000-pt subsamples, seeds 100-104
        |              --> emits (mu, sigma) natively
        |
        +--> QRNN x5   arch x feature subset x width
        |              --> 9 deciles --> (mu, sigma) by probit fit
        |
        +--> QRGB x5   XGBoost, varied depth/child weight/subsample/colsample
        |              --> 9 deciles --> (mu, sigma) by probit fit
        |
        v        fit on T1, predict on T2+T3+T4
   15 Gaussians  F_i = N(mu_i, sigma_i^2)
        |
        v        weights w solved on T3
   mixture  F = sum_i w_i F_i          w_i = 0 drops model i
        |
        v        scored on T4
   invert F at {0.05, 0.50, 0.95}  -->  (3, n_test)
\end{codetiny}

\noindent
\textbf{Probit fit.} If $y_q = \mu + \Phi^{-1}(q)\sigma$, an OLS of the nine deciles $y_q$ on
$\Phi^{-1}(q)$ returns intercept $\mu$ and slope $\sigma$.

\medskip
\noindent
\textbf{Weights.} For a Gaussian mixture the CRPS is closed form (Baringhaus--Franz),
\begin{equation}
  \mathrm{CRPS}(F,y) = \mathbb{E}|Y-y| - \tfrac12 \mathbb{E}|Y-Y'|
  = \sum_{i,k} \alpha_{ik} w_i w_k + \sum_i \beta_i w_i ,
\end{equation}
so the weights are a small QP on the simplex,
\begin{equation}
  \min_{w}\; w^{\top} Q w + c^{\top} w
  \quad\text{s.t.}\quad \mathbf{1}^{\top} w = 1,\; w \geq 0 .
\end{equation}
$Q$ is indefinite ($\alpha_{ii} = -\sigma_i/\sqrt{\pi}$); eliminating the equality constraint
gives a reduced matrix that is checked for positive definiteness on every run rather than
assumed. Solved by multi-start SLSQP.

\medskip
\noindent
Selection and weighting happen in one solve --- $w_i = 0$ drops model $i$. Reported alongside:
MAPE-guided weights, simple average, best single model on T3.

\subsection{Model B --- iQRNN}

\begin{code}
input (103)
  Linear(103 -> 500) -> PReLU -> Dropout(p) -> GaussianNoise(0.3)
  Linear(500 -> 500) -> PReLU -> Dropout(p) -> GaussianNoise(0.3)
  Linear(500 ->   3)                     output width = the quantile set

  Adam lr 1e-3, batch 168
  bad = epochs since best val pinball
      bad >= 20         stop
      every 5 bad       lr *= 0.4, floor 1e-5
  best weights restored;  p = 0.5 and p = 0 both run, 5 seeds each
\end{code}

\noindent
\textbf{Loss.} The output width \emph{is} the quantile set and the objective \emph{is} the
reported metric, so one fit produces every quantile:
\begin{equation}
  \mathcal{L} = \sum_{q \in Q} \max\!\big\{ q(y - \yhat_q),\; (q-1)(y - \yhat_q) \big\},
  \qquad Q = \{0.05,\,0.50,\,0.95\} .
\end{equation}
Outputs are sorted to remove crossing. Noise injection is the largest single ablation effect
in the paper.

% ============================================================================
\section{Metrics and results}

\noindent
$\alpha = 0.10$, interval $[\ell, u] = [\yhat_{0.05}, \yhat_{0.95}]$, averaged over
$N$ days $\times$ $H$ steps:
\begin{align}
  \mathrm{Pinball} &= \tfrac{1}{|Q| N H} \sum_{q}\sum_{d,h}
      \max\{q(y-\yhat_q),\,(q-1)(y-\yhat_q)\} \\
  \mathrm{PICP} &= \tfrac{1}{NH}\sum_{d,h} \ind\{\ell \leq y \leq u\},
  \qquad
  \mathrm{MPIW} = \tfrac{1}{NH}\sum_{d,h} (u - \ell) \\
  \mathrm{Winkler} &= \tfrac{1}{NH}\sum_{d,h}\Big[(u-\ell)
      + \tfrac{2}{\alpha}(\ell-y)\ind\{y<\ell\}
      + \tfrac{2}{\alpha}(y-u)\ind\{y>u\}\Big]
\end{align}
MAE / RMSE / MAPE use the median $\yhat_{0.50}$.

\medskip
\noindent
\textbf{Units.} The target is z-scored with the training days' mean and sd; everything is
reported in those units except MAPE, which is mapped back to physical units because the
z-scored target crosses zero.

\begin{center}
\small
\begin{tabular}{@{}l l l c c c c c c@{}}
\toprule
 & model & Pinball & PICP & MPIW & Winkler & MAE & RMSE & MAPE \\
\midrule
\texttt{spain} & A CRPS-LOP & 0.0867 $\pm$ 0.0002 & 0.8826 & 1.2348 & 2.1218 & 0.3079 & 0.4903 & 4.82 \\
               & B iQRNN-500 & 0.1008 $\pm$ 0.0006 & 0.8441 & 1.3587 & 2.3612 & 0.3687 & 0.5325 & 5.83 \\
\addlinespace
\texttt{ercot} & A CRPS-LOP & 0.1154 $\pm$ 0.0004 & 0.8598 & 1.6311 & 2.2960 & 0.4625 & 0.5710 & 18.92 \\
               & B iQRNN-500 & 0.0979 $\pm$ 0.0010 & 0.8358 & 1.3673 & 1.9799 & 0.3895 & 0.4863 & 15.32 \\
\addlinespace
\texttt{elia}  & A CRPS-LOP & 0.0573 $\pm$ 0.0002 & 0.8949 & 0.8701 & 1.2977 & 0.2139 & 0.3025 & 3.35 \\
               & B iQRNN-500 & 0.0433 $\pm$ 0.0003 & 0.9035 & 0.6930 & 0.9292 & 0.1670 & 0.2217 & 2.61 \\
\bottomrule
\end{tabular}
\end{center}

\noindent
Three caveats before quoting these.

\begin{enumerate}[leftmargin=1.6em,itemsep=2pt]
  \item \textbf{No consistent winner.} B wins on \texttt{ercot} and \texttt{elia}, loses on
        \texttt{spain}. The \texttt{spain} loss coincides with B stopping at epoch 3--5 there
        and on \texttt{ercot} (\texttt{elia} trains 58--61 epochs), so it is likely
        non-convergence rather than a property of the method.
  \item \textbf{A's combination loses to its own best member on \texttt{ercot}} (CRPS 0.3234
        vs 0.3143), while winning $\sim$4\% elsewhere. \texttt{ercot}'s T3 is 71 days of high
        summer and T4 runs into January --- weights fitted in one season, applied across a
        change of season. A consequence of a two-year record, not of the combination step.
  \item \textbf{\texttt{ercot}'s MAPE is not comparable.} Its target is \emph{net} load, whose
        test minimum is $0.353\times$ its median. Its z-unit MAE is $1.5\times$
        \texttt{spain}'s while its MAPE is $3.9\times$; the gap is the artefact.
\end{enumerate}

\noindent
Also worth stating: the three datasets do not pose the same task. \texttt{elia}'s
\texttt{known\_cols} contains the official TSO day-ahead load forecast; the other two do not.
Part of \texttt{elia}'s advantage is the TSO's. And the feature sets are asymmetric by design
--- A gets no exogenous columns (Li's specification), B does --- so on \texttt{elia} B sees
that forecast and A does not.

\end{document}
